\chapter{CDP, LLDP and Documenting a Network}\label{ch:discovery}
\bp{2.3}

\begin{outcomes}
By the end of this chapter you will be able to:
\begin{tight}
  \item \textbf{Explain} what discovery protocols advertise and what they cannot
        tell you.
  \item \textbf{Configure} CDP and LLDP, including disabling them where they
        should not run.
  \item \textbf{Validate} a network diagram against the devices themselves, which
        is precisely what topic \tp{2.3} asks.
  \item \textbf{Use} discovery output to find a mismatch --- native VLAN, duplex
        or an undocumented device.
\end{tight}
\end{outcomes}

\begin{prereq}
\chref{trunking} for native VLANs, \chref{physical} for duplex, \chref{cli} for
filtering output.
\end{prereq}

\begin{keyterms}
CDP $\bullet$ LLDP $\bullet$ LLDP-MED $\bullet$ neighbour $\bullet$ holdtime
$\bullet$ advertisement $\bullet$ capability code $\bullet$ platform $\bullet$
local interface $\bullet$ port ID
\end{keyterms}

\begin{examnote}[The verb here is \emph{validate}]
Topic \tp{2.3} is ``validate the accuracy of network documentation using CDP and
LLDP''. It is not asking you to describe the protocols. It is asking you to take
a diagram that claims \cmd{SW1 Gi0/1} connects to \cmd{R1 Gi0/0}, and prove from
the devices whether that is true. Expect to be shown a diagram and some
\cmd{show cdp neighbors} output that contradicts it.
\end{examnote}

\section{Two protocols, one job}

\begin{longtable}{@{}L{3.0cm}L{5.6cm}L{6.4cm}@{}}
\toprule
\rowcolor{primary!15}
\textbf{} & \textbf{CDP} & \textbf{LLDP} \\
\midrule
\endhead
Standard & Cisco proprietary & IEEE 802.1AB, open \\
\rowcolor{lightbg}
Default state & On (most Cisco platforms) & Off (most Cisco platforms) \\
Timers & Advertise 60~s, hold 180~s & Advertise 30~s, hold 120~s \\
\rowcolor{lightbg}
Sees & Cisco devices & Anything that speaks LLDP, including other vendors,
printers and phones \\
Voice extension & --- & LLDP-MED, used to tell phones their voice VLAN \\
\bottomrule
\end{longtable}

Run both in a mixed estate: CDP for the Cisco gear, LLDP for everything else.

\begin{definitionbox}[What a neighbour advertisement carries]
Device identifier, the local and remote interface, the platform and capabilities,
and depending on the protocol and platform: IP address, IOS version, native VLAN,
duplex and power requirements.

Crucially, a neighbour is \textbf{directly connected at layer 2}. A device two
switches away never appears. That is exactly what makes the output usable as
evidence of physical cabling.
\end{definitionbox}

\section{Reading the output}

\begin{verify}{SW1}{show cdp neighbors}
Capability Codes: R - Router, T - Trans Bridge, B - Source Route Bridge
                  S - Switch, H - Host, I - IGMP, r - Repeater, P - Phone

Device ID    Local Intrfce   Holdtme   Capability   Platform    Port ID
R1           Gig 0/1          142        R S I      C1111       Gig 0/0/0
SW2          Gig 0/2          167        S I        WS-C2960    Gig 0/1
SW3          Gig 0/3          121        S I        WS-C2960    Gig 0/24
AP-EAST      Gig 0/10         155        H          AIR-AP2802  Gig 0
\end{verify}

Read it as four columns that matter: \textbf{Device ID} (who), \textbf{Local
Intrfce} (my port), \textbf{Port ID} (their port), \textbf{Capability} (what kind
of device). Those four are the cabling diagram, taken from the devices rather
than from a drawing somebody updated in 2019.

\begin{verify}{SW1}{show cdp neighbors detail}
-------------------------
Device ID: R1
Entry address(es):
  IP address: 10.0.0.1
Platform: Cisco C1111-4P,  Capabilities: Router Switch IGMP
Interface: GigabitEthernet0/1,  Port ID (outgoing port): GigabitEthernet0/0/0
Holdtime : 142 sec

Version :
Cisco IOS XE Software, Version 17.09.04a

Duplex: full
Native VLAN: 999
\end{verify}

\cmd{detail} adds the IP address, the software version, the duplex and the native
VLAN --- which is how a native VLAN mismatch gets caught before it causes a
security incident.

\begin{config}{running both protocols, and switching them off where they do not belong}
! CDP is on by default; LLDP is not.
lldp run
!
! Turn CDP off on a port facing anything you do not control.
! Advertising your platform and IOS version to a guest network,
! an ISP handoff or an untrusted uplink is free reconnaissance.
interface GigabitEthernet0/24
 description ISP handoff - no discovery
 no cdp enable
 no lldp transmit
 no lldp receive
!
! Globally, if policy demands it:
!   no cdp run
\end{config}

\begin{alertbox}[Do not disable it everywhere on principle]
There is a reflex to switch CDP off across the estate ``for security''. Weigh it
properly: internally, CDP and LLDP are how you validate cabling, how phones learn
their voice VLAN (\chref{edgehost}), and how native VLAN and duplex mismatches
are detected. Disable them on untrusted edges --- ISP links, guest ports, DMZ
segments --- and leave them running inside.
\end{alertbox}

\section{Validating a diagram}

\begin{worked}[The diagram says one thing, the switch says another]
The documentation claims:

\begin{tight}
  \item \cmd{SW1 Gi0/2} $\rightarrow$ \cmd{SW2 Gi0/1}
  \item \cmd{SW1 Gi0/3} $\rightarrow$ \cmd{SW3 Gi0/1}
\end{tight}

The output above shows \cmd{Gi0/2} to \cmd{SW2 Gi0/1} --- correct --- but
\cmd{Gi0/3} to \cmd{SW3 \textbf{Gig 0/24}}, not \cmd{Gi0/1}.

\textbf{What that means in practice.} Somebody re-patched \cmd{SW3} and did not
update the diagram. It may be harmless. It may not: if \cmd{SW3 Gi0/24} is
configured as an access port while \cmd{Gi0/1} was the trunk, the uplink is
carrying one VLAN instead of all of them, and the fault will surface later as
``the new VLAN does not work on the third floor''.

\textbf{What to do.} Correct the diagram, then check that the port actually in use
carries the configuration the design intended. Discovery output tells you what is
plugged in; it does not tell you whether the port is configured correctly. That
is a second check.
\end{worked}

\begin{pitfall}
\begin{tight}
  \item \textbf{Expecting to see non-adjacent devices.} Neighbours are one hop.
        A device that does not appear may simply be further away.
  \item \textbf{Concluding a device is absent when CDP is off.} A non-Cisco
        switch, or one with CDP disabled, is invisible to CDP and may still be
        there. Cross-check with LLDP and the MAC address table.
  \item \textbf{Trusting stale entries.} Holdtime counts down from 180~s. An
        entry can persist for three minutes after a device is unplugged.
  \item \textbf{Forgetting \cmd{lldp run}.} Without it, \cmd{show lldp neighbors}
        is empty and you may wrongly conclude nothing is connected.
  \item \textbf{Reading the port columns backwards.} \cmd{Local Intrfce} is
        yours; \cmd{Port ID} is theirs.
\end{tight}
\end{pitfall}

\begin{breakfix}[An undocumented switch appears in the discovery table.]
An audit of \cmd{SW1} shows a neighbour that is not on any diagram:

\begin{verify}{SW1}{show cdp neighbors GigabitEthernet0/14 detail}
-------------------------
Device ID: Switch
Entry address(es):
Platform: cisco WS-C2960-8TC-L,  Capabilities: Switch IGMP
Interface: GigabitEthernet0/14,  Port ID (outgoing port): FastEthernet0/1
Holdtime : 131 sec

Version :
Cisco IOS Software, C2960 Software, Version 12.2(55)SE
\end{verify}

\textbf{Diagnosis.} An eight-port switch with the default hostname \cmd{Switch},
no management address, and a decade-old IOS version, connected to an access port
--- somebody has plugged a desk switch into a wall socket to get more ports. It is
unmanaged, it participates in spanning tree, and it can influence the root
election (\chref{stp}).

\textbf{Fix.} Short term: confirm what is downstream of it before unplugging
anything, because there are probably users behind it. Long term, this is exactly
what \cmd{spanning-tree bpduguard} on access ports is for --- with it configured,
the port would have errdisabled the moment the switch was connected and the
problem would have announced itself instead of being found in an audit months
later.
\end{breakfix}

\begin{lab}{Draw the network from the devices}{Packet Tracer}
\textbf{Topology.} Four switches and a router, cabled by your instructor or a
partner without telling you how, plus a deliberately wrong diagram.

\begin{tasks}
  \item Using only \cmd{show cdp neighbors} from each device, draw the physical
        topology. Do not look at the cables.
  \item Compare your drawing with the supplied diagram and list every
        discrepancy.
  \item Run \cmd{show cdp neighbors detail} on one link and record the IOS
        version, duplex and native VLAN of the neighbour.
  \item Create a native VLAN mismatch on one trunk and find the resulting log
        message.
  \item Enable \cmd{lldp run} everywhere and compare \cmd{show lldp neighbors}
        with the CDP view. Account for any difference.
  \item Disable CDP on one port and confirm the neighbour disappears --- noting
        how long it takes.
  \item \textbf{Extension.} Write a short procedure a junior technician could
        follow to validate a patching change using discovery output alone.
\end{tasks}
\end{lab}

\begin{checkpoint}
\begin{tightnum}
  \item \cmd{show cdp neighbors} lists a device you know is three switches away.
        What must be true?
  \item You unplug a device and it still appears in the table. Why, and for how
        long?
  \item Which single command would tell you a neighbour's IOS version and native
        VLAN?
\end{tightnum}
\end{checkpoint}

\begin{chaptersummary}
\begin{tight}
  \item CDP is Cisco and on by default; LLDP is open and off by default. Run both
        in a mixed estate.
  \item Neighbours are directly connected only --- which is what makes the output
        evidence of cabling.
  \item \cmd{Local Intrfce} is your port, \cmd{Port ID} is theirs. Those two
        columns plus the device name are the diagram.
  \item \cmd{detail} adds IP, version, duplex and native VLAN, and catches
        mismatches before they cause incidents.
  \item Disable discovery on untrusted edges, not everywhere. Holdtime means
        entries persist for minutes after a device is removed.
\end{tight}
\end{chaptersummary}

\begin{reviewq}
  \item \textbf{(MCQ)} Which command enables LLDP globally?
        (a)~\cmd{lldp enable} (b)~\cmd{lldp run} (c)~\cmd{lldp transmit}
        (d)~\cmd{no lldp disable}
  \item \textbf{(MCQ)} In \cmd{show cdp neighbors}, which column names the port
        on the \emph{remote} device?
        (a)~Local Intrfce (b)~Port ID (c)~Platform (d)~Capability
  \item \textbf{(Short)} Give two reasons to leave CDP running inside a network
        and one reason to disable it at an edge.
  \item \textbf{(Short)} A diagram says two switches are connected. Neither
        appears in the other's CDP table. Give three explanations.
  \item \textbf{(Applied)} You are handed a network with no documentation.
        Describe a repeatable procedure, naming the commands, to produce an
        accurate physical topology diagram and state two things your procedure
        cannot establish.
\end{reviewq}
